%%%%%%%%%%%%%%%%%%%%%%%%%%%%%%%%%%%%%%%%%%%%%%%%%%%%%%%%%%%%%%%%%%%%%%%%%%%
% Función: guía de los formularios de TFG implementados para la EPS-UAH.
%
% Uso: consulta el PDF para preparar los formularios. Normalmente no debes
% modificar este fichero, salvo para mantener la documentación de la plantilla.
%%%%%%%%%%%%%%%%%%%%%%%%%%%%%%%%%%%%%%%%%%%%%%%%%%%%%%%%%%%%%%%%%%%%%%%%%%%

\documentclass[10pt,a4paper]{article}
\usepackage[T1]{fontenc}
\usepackage[utf8]{inputenc}
\usepackage{lmodern}
\usepackage[spanish,es-lcroman]{babel}
\usepackage[margin=2cm]{geometry}
\usepackage{enumitem,tabularx,booktabs}
\usepackage[colorlinks=true,urlcolor=blue,linkcolor=blue]{hyperref}
\usepackage{xurl}
\usepackage{microtype}
\usepackage{xcolor,soul}
\sethlcolor{yellow}
\setcounter{tocdepth}{1}
\setlength{\parindent}{0pt}
\setlength{\parskip}{4pt}
\setlength{\emergencystretch}{2em}
\setlist[itemize]{leftmargin=*,nosep,topsep=3pt}
\newcommand{\variable}[1]{\texttt{\textbackslash{}}\nolinkurl{#1}}
\newcommand{\fichero}[1]{\path{#1}}
\title{Guía del papeleo disponible para los TFGs en la EPS-UAH\\\large Escuela Politécnica Superior --- Universidad de Alcalá}
\author{}
\date{}

\begin{document}
\begin{center}
{\LARGE Guía del papeleo disponible para los TFGs en la EPS-UAH\par}
\medskip
{\large Escuela Politécnica Superior --- Universidad de Alcalá\par}
\end{center}

{\small
\setlength{\parskip}{0pt}
\tableofcontents
}
\medskip

\section{Qué encontrarás aquí}

Hemos preparado seis formularios específicos de la EPS-UAH y un informe de control de plagio basado en la normativa general de TFG de la UAH. Todos reutilizan los datos de tu trabajo para que no tengas que introducirlos una y otra vez. En esta guía te contamos para qué sirve cada uno, qué variables utiliza y qué tendrás que revisar antes de presentarlo. Los formularios específicos de la EPS-UAH están actualizados a su normativa de TFG de octubre de 2026. El informe general corresponde al anexo II de la normativa UAH y puede ser útil fuera de la EPS-UAH. Esta guía no sustituye los formularios y procedimientos oficiales vigentes de cada centro.

Puedes localizar los documentos mediante este resumen. Cada nombre enlaza con su explicación dentro de la guía.

\medskip
\begin{tabularx}{\textwidth}{@{}Xl@{}}
\toprule
Documento & Anexo de la normativa \\
\midrule
\hyperref[form:conjunta]{Autorización conjunta de publicación} & VII (EPS-UAH) \\
\hyperref[form:extranjero]{Autorización del tutor extranjero} & VIII (EPS-UAH) \\
\hyperref[form:confidencialidad]{Compromisos de confidencialidad} & VI (EPS-UAH) \\
\hyperref[form:convenio]{Convenio de cooperación educativa} & V (EPS-UAH) \\
\hyperref[form:informe]{Informe del tutor y rúbrica} & III (EPS-UAH) \\
\hyperref[form:plagio]{Informe general de control de plagio} & II (UAH) \\
\hyperref[form:acta]{Acta de defensa y rúbrica del tribunal} & IX (EPS-UAH) \\
\bottomrule
\end{tabularx}

\section{Antes de generar los documentos}

La marca \hl{N/D} significa \emph{no disponible} y aparece resaltada en amarillo en los campos sin información. Revisa cuáles debes completar antes de presentar los formularios.

Abre \fichero{Config/myconfig.tex} y sustituye los datos de ejemplo por los tuyos. Selecciona tu titulación de grado mediante \variable{myDegree} y el idioma mediante \variable{myLanguage}. Para los formularios específicos de la EPS-UAH, selecciona una titulación de ese centro. El identificador del grado determina la universidad, el centro y el tipo de trabajo. La plantilla completa automáticamente esos datos para cada formulario.

Tu nombre completo, \variable{myAuthorFullName}, se obtiene normalmente de \variable{myAuthorName} y \variable{myAuthorSurname}. Cuando una ficha mencione el nombre completo, revisa estas dos variables. Define \variable{myBookTitleSpanish} y, cuando proceda, \variable{myBookTitleEnglish}: la plantilla selecciona el título según \variable{myLanguage}. Las fichas siguientes indican los datos particulares de cada formulario, incluidos los utilizados por sus fragmentos auxiliares.

La fecha \variable{myPaperworkDate} toma por defecto el valor de \variable{myThesisDepositDate}; la versión inglesa hace lo mismo con \variable{myThesisDepositDateEnglish}. Puedes modificar esas fechas sin cambiar los comandos del formulario. No confundas estas variables con la fecha de firma digital o con la fecha del equipo, utilizadas en algunos documentos.

Configura el departamento donde realizas el trabajo mediante \variable{myDepartment} y su versión inglesa mediante \variable{myDepartmentEnglish}. Son datos independientes de las afiliaciones del tutor y del cotutor: si los dejas vacíos, permanecen vacíos. Los formularios administrativos en español utilizan el primer campo. Configura por separado los departamentos e instituciones de tus tutores; los campos que solicitan su afiliación muestran ambos datos cuando están disponibles.

\subsection{Género gramatical y singular o plural}

La plantilla adapta expresiones como \emph{autor/autora}, \emph{tutor/tutora}, \emph{el/la} o \emph{D./Dª.} mediante las variables de género. Debes escribir exactamente \texttt{male} o \texttt{female}; el nombre de la persona no se utiliza para deducirlo y la implementación no valida otros valores. Comprueba especialmente \variable{myAuthorGender}, \variable{myAcademicTutorGender} y \variable{myCoTutorGender}. Las fichas te indican cuáles afectan realmente a cada documento.

Si no tienes cotutor, deja \variable{myCoTutorFullName} vacío: \verb|\newcommand{\myCoTutorFullName}{}|. Su presencia controla las secciones y firmas opcionales y, donde se usan los comandos correspondientes, el singular o plural. No escribas \hl{N/D} como nombre para indicar su ausencia, porque entonces la plantilla lo tratará como una persona configurada.

Por ejemplo, la autorización conjunta puede mostrar \emph{la autora} y \emph{la tutora} si ambas variables de género son \texttt{female}. En los compromisos individuales se obtiene \emph{El Profesor} o \emph{La Profesora}; el tutor \emph{tutoriza}, el cotutor \emph{cotutoriza} y un miembro del tribunal \emph{evaluará}. No todos los textos se adaptan: algunos rótulos institucionales son fijos. Los comandos comunes de plural para dos supervisores utilizan actualmente \emph{tutores/directores}, también si ambas personas son mujeres.

\subsection{Casillas configurables y casillas fijas}

No todas las casillas se controlan desde \fichero{Config/myconfig.tex}. Las fichas siguientes distinguen las que dependen de una variable de las que están escritas directamente en el formulario. Respeta las mayúsculas de los valores: \texttt{Yes}/\texttt{No} para publicación y \texttt{YES}/\texttt{NO} para la propuesta de Matrícula de Honor. Tras cambiar la configuración, vuelve a compilar y comprueba las marcas en el PDF.

Si una casilla es fija, debes editar el fichero \texttt{.tex} del formulario cuando su selección no corresponda a tu caso: \verb|$\XBox$| imprime una casilla marcada y \verb|$\Box$| una vacía. Intercambia ambos comandos para cambiar la respuesta y deja marcada solo la opción que proceda. Cambiar una variable de publicación no modifica las casillas fijas.

\section{Publicación y confidencialidad}

\subsection{Autorización conjunta de publicación}\label{form:conjunta}

\fichero{TFG-AnexoVII-AutorizacionATodosPubAbierto-EPS-UAH.tex}. Anexo VII de la normativa EPS-UAH de octubre de 2026. Reúne la autorización y las firmas del estudiante, el tutor y, si existe, el cotutor.

Configura los siguientes datos:
\begin{itemize}
  \item Estudiante y título: \variable{myAuthorFullName}, \variable{myAuthorDNI}, \variable{myBookTitleSpanish} y \variable{myDegree} (para la titulación).
  \item Tutor: \variable{myAcademicTutorFullName}, \variable{myAcademicTutorDNI}, \variable{myAcademicTutorDepartment} y \variable{myAcademicTutorInstitution}.
  \item Cotutor, si existe: \variable{myCoTutorFullName}, \variable{myCoTutorDNI}, \variable{myAcademicCotutorDepartment} y \variable{myAcademicCotutorInstitution}.
  \item Publicación: \variable{myAuthorizationOpenPublishing}, con \texttt{Yes} o \texttt{No}. Este formulario no muestra opciones de embargo.
  \item Género: \variable{myAuthorGender}, \variable{myAcademicTutorGender} y \variable{myCoTutorGender}; adaptan los rótulos y las firmas. La firma del cotutor desaparece si su nombre está vacío.
\end{itemize}
Las casillas dependen de \variable{myAuthorizationOpenPublishing}: \texttt{Yes} marca \emph{AUTORIZAN} y deja vacía \emph{NO AUTORIZAN}; \texttt{No} hace lo contrario. \variable{myAuthorizationOpenPublishingEmbargoMonths} no interviene en este formulario. Revisa la decisión de autorización y recoge las firmas correspondientes.

\subsection{Autorización del tutor extranjero}\label{form:extranjero}

\fichero{TFG-AnexoVIII-AutorizacionTutorExtranjeroPubAbierto-EPS-UAH.tex}. Anexo VIII de la normativa EPS-UAH de octubre de 2026. Está redactado en inglés y utiliza los datos del \emph{tutor académico}, no los del cotutor: comprueba que el nombre y los campos de universidad y ciudad extranjeras describan a esa misma persona.

Configura \variable{myAcademicTutorFullName}, \variable{myAcademicTutorForeignUniversity}, \variable{myAcademicTutorForeignCity}, \variable{myAuthorFullName}, \variable{myBookTitleEnglish} y \variable{myPaperworkDateEnglish}. El texto utiliza \emph{Professor} y no emplea variables de género. Las casillas son fijas: \emph{I agree} está marcada y \emph{I disagree} está vacía; no dependen de \variable{myAuthorizationOpenPublishing}. Si el tutor no autoriza, debes editar este fichero e intercambiar \verb|$\XBox$| y \verb|$\Box$| en esas dos opciones. Comprueba la decisión y obtén la firma y el sello indicados.

\subsection{Compromisos de confidencialidad}\label{form:confidencialidad}

\fichero{TFG-AnexoVI-CompromisoConfidencialidad-EPS-UAH.tex}. Anexo VI de la normativa EPS-UAH de octubre de 2026. Genera un compromiso para el tutor y otro para cada persona opcional cuyo nombre esté configurado: cotutor, presidente, dos vocales y suplente. El resultado es un único PDF con hasta seis compromisos; en cada uno firman la persona correspondiente y el estudiante.

Revisa estas variables:
\begin{itemize}
  \item Estudiante y trabajo: \variable{myAuthorFullName}, \variable{myAuthorDNI}, \variable{myAuthorGender}, \variable{myBookTitleSpanish} o \variable{myBookTitleEnglish} y \variable{myPaperworkDate}.
  \item Tutor: \variable{myAcademicTutorFullName}, \variable{myAcademicTutorDNI} y \variable{myAcademicTutorGender}.
  \item Cotutor: \variable{myCoTutorFullName}, \variable{myCoTutorDNI} y \variable{myCoTutorGender}.
  \item Presidente: \variable{myTribunalPresident}, \variable{myTribunalPresidentDNI} y \variable{myTribunalPresidentGender}.
  \item Vocal primero: \variable{myTribunalFirstSpokesperson}, \variable{myTribunalFirstSpokespersonDNI} y \variable{myTribunalFirstSpokespersonGender}.
  \item Vocal segundo: \variable{myTribunalSecondSpokesperson}, \variable{myTribunalSecondSpokespersonDNI} y \variable{myTribunalSecondSpokespersonGender}.
  \item Suplente: \variable{myTribunalAlternateMember}, \variable{myTribunalAlternateMemberDNI} y \variable{myTribunalAlternateMemberGender}.
\end{itemize}
Un nombre opcional vacío omite esa persona y un DNI de profesor vacío deja una línea para completarlo. La fecha procede de \variable{myPaperworkDate}; el título procede de \variable{myBookTitleSpanish} o \variable{myBookTitleEnglish}, según \variable{myLanguage}. El formulario no se activa mediante \variable{myConfidentialContent}: genéralo únicamente si lo necesitas.

\section{Cooperación educativa}

\subsection{Anexo al convenio}\label{form:convenio}

\fichero{TFG-AnexoV-ConvenioCooperacion-EPS-UAH.tex}. Anexo V de la normativa EPS-UAH de octubre de 2026. Describe el trabajo que se realizará en una entidad externa. La persona responsable en esa entidad es el cotutor externo, de acuerdo con el artículo 3.3.b de la normativa. La plantilla utiliza esta denominación también en el encabezado de contacto y en los términos del acuerdo para mantenerla coherente con el bloque de firmas.

Utiliza los siguientes datos:
\begin{itemize}
  \item Estudiante: \variable{myAuthorFullName}, \variable{myAuthorDNI}, \variable{myAuthorTelephone} y \variable{myAuthorEmail}.
  \item Tutor académico: \variable{myAcademicTutorFullName}, \variable{myAcademicTutorDNI}, \variable{myAcademicTutorDepartment}, \variable{myAcademicTutorInstitution} y \variable{myAcademicTutorEmail}.
  \item Entidad y cotutor externo: \variable{myCoTutorFullName}, \variable{myAcademicCotutorInstitution}, \variable{myCoTutorPosition}, \variable{myCoTutorDNI} y \variable{myCoTutorEmail}. Sin cotutor, los campos derivados de él muestran \hl{N/D}.
  \item Firma por la UAH: \variable{mySchoolDeputyDirectorFullName} y \variable{mySchoolDeputyDirectorGender}, con \texttt{male} o \texttt{female}. El nombre aparece sobre el cargo \emph{Subdirector de la} o \emph{Subdirectora de la}, seguido del nombre de la escuela en la línea siguiente. Si el nombre está vacío, su línea permanece en blanco.
  \item Trabajo: \variable{myDegree}, \variable{myBookTitleSpanish} y \variable{myThesisAcademicYear}. La titulación seleccionada proporciona el centro y el tipo de trabajo.
  \item Género: \variable{myAuthorGender} adapta el tratamiento del estudiante. \variable{myAcademicTutorGender} adapta el rótulo del tutor o tutora de la UAH. \variable{myCoTutorGender} adapta el encabezado de contacto, las dos cláusulas del acuerdo y la firma del cotutor externo, con las formas \emph{Cotutor/Cotutora} y \emph{El cotutor/La cotutora}.
\end{itemize}
Completa los campos todavía marcados como \hl{N/D}: identificación y dirección de la entidad, teléfonos, fechas de inicio y fin, horario y resumen del trabajo. No tienen variables propias en \fichero{myconfig.tex}; puedes completar el formulario o adaptar sus campos en el fuente. La fecha impresa corresponde a la fecha del equipo durante la compilación, no \variable{myPaperworkDate}. Revisa también las tres firmas: subdirección de la escuela, cotutor externo y estudiante.

\section{Informes y evaluación}

\subsection{Informe del tutor y rúbrica}\label{form:informe}

\fichero{TFG-AnexoIII-InformeTutorYRubrica-EPS-UAH.tex}. Anexo III de la normativa EPS-UAH de octubre de 2026. Incluye el informe de plagio, los datos y firmas del tutor y del cotutor, si existe, y la rúbrica del tutor.

La primera página utiliza \variable{myAuthorFullName}, \variable{myAuthorEmail}, \variable{myDegree}, \variable{myBookTitleSpanish} y, con \variable{myLanguage} igual a \texttt{english}, \variable{myBookTitleEnglish}. El acta de evaluación añade \variable{myAcademicTutorFullName}, \variable{myAcademicTutorDepartment}, \variable{myCoTutorFullName}, \variable{myAcademicCotutorDepartment}, \variable{myAcademicCotutorInstitution}, \variable{myThesisAcademicYear} y \variable{myAcademicTutorGrade}, de cero a diez. Los rótulos y firmas dependen de \variable{myAcademicTutorGender} y \variable{myCoTutorGender}.

Completa los resultados del control de plagio, las observaciones, la conclusión, los comentarios del tutor y la rúbrica. La fecha indicada es la de la firma digital; \variable{myThesisDefenseDate} no se imprime actualmente en este documento.

\subsection{Informe general UAH de control de plagio}\label{form:plagio}

\fichero{TFG-AnexoII-InformeTutorPlagio-UAH.tex}. Está basado en el anexo II de la normativa general de TFG de la UAH, de febrero de 2024. Es el único formulario administrativo definido en esa normativa general y puede ser útil en otros centros de la UAH, no solo en la EPS-UAH. No incorpora la rúbrica de evaluación de la EPS-UAH. Comprueba los requisitos concretos de tu centro antes de presentarlo.

Utiliza \variable{myAuthorFullName}, \variable{myAuthorDNI}, \variable{myAuthorEmail}, \variable{myDepartment}, \variable{myDegree}, \variable{myBookTitleSpanish} y, con \variable{myLanguage} igual a \texttt{english}, \variable{myBookTitleEnglish}. Revisa también \variable{myAcademicTutorFullName}, \variable{myAcademicTutorPosition}, \variable{myAcademicTutorEmail}, \variable{myAcademicTutorDepartment}, \variable{myCoTutorFullName}, \variable{myCoTutorPosition}, \variable{myCoTutorEmail}, \variable{myAcademicCotutorDepartment} y \variable{myAcademicCotutorInstitution}. Los rótulos del tutor y cotutor se adaptan con \variable{myAcademicTutorGender} y \variable{myCoTutorGender}; no se emplea \variable{myAuthorGender} para los rótulos del estudiante de este informe.

Completa coincidencias, observaciones y conclusión del análisis y recoge las firmas. La fecha es la de la firma digital, no \variable{myThesisDefenseDate}.

\subsection{Acta de defensa y rúbrica del tribunal}\label{form:acta}

\fichero{TFG-AnexoIX-ActaYRubricaDefensa-EPS-UAH.tex}. Anexo IX de la normativa EPS-UAH de octubre de 2026. Incluye las calificaciones y observaciones del tribunal y su rúbrica.

Configura \variable{myAuthorFullName}, \variable{myAuthorEmail}, \variable{myDegree}, \variable{myBookTitleSpanish}, \variable{myBookTitleEnglish} (si \variable{myLanguage} es \texttt{english}), \variable{myAcademicTutorFullName}, \variable{myAcademicTutorDepartment}, \variable{myCoTutorFullName} y \variable{myThesisAcademicYear}. El departamento solicitado aquí es el del tutor, no el departamento independiente del trabajo. Para las notas utiliza \variable{myAcademicTutorGrade} y \variable{myTribunalGrade}, ambas de cero a diez, sin ponderarlas previamente. La nota final es \(0.7 \times \mbox{nota del tutor} + 0.3 \times \mbox{nota del tribunal}\). Por ejemplo, con 8 y 6 obtendrás 7.4. La plantilla rechaza notas fuera de ese intervalo y detiene la compilación si alguna parte está suspensa pero la media daría aprobado, de acuerdo con el artículo 9.3. En ese caso, consulta al departamento cómo consignar la nota: no inventamos una calificación máxima para el suspenso.

Los rótulos de los tutores se adaptan mediante \variable{myAcademicTutorGender} y \variable{myCoTutorGender}. La rúbrica utiliza \variable{myThesisDefenseDate}, \variable{myTribunalPresident} y \variable{myTribunalFirstSpokesperson}, con sus variables de género. Si procede una tercera firma, completa \variable{myTribunalSecondSpokesperson} y \variable{myTribunalSecondSpokespersonGender}: aparecerá como \emph{Secretario / Vocal 2º} o \emph{Secretaria / Vocal 2ª}. Si el tribunal tiene solo dos titulares, deja vacío el nombre del segundo vocal y se omite esa tercera firma. El suplente no firma salvo que sustituya a un titular. Completa las observaciones y comprueba la composición del tribunal antes de presentar el documento.

Las casillas de propuesta para Matrícula de Honor dependen de \variable{myTribunalMHProposal}: \texttt{YES} marca \emph{Sí} y deja vacía \emph{No}; \texttt{NO} hace lo contrario. La propuesta no se deduce de las notas ni de su media ponderada: debes configurar la decisión del tribunal y comprobar que cumple los requisitos aplicables.

\medskip
\begin{minipage}{\linewidth}
\setlength{\parskip}{4pt}
\section{Cómo generar los documentos}

No necesitas \texttt{make}. Abre el formulario que necesites en tu editor de \LaTeX{}, mantén la estructura de directorios y compílalo con \texttt{pdflatex} desde \fichero{PapeleoTFG/}. Si tu editor lo requiere, establece esa carpeta como directorio de trabajo. Estas plantillas no necesitan \texttt{biber} ni \texttt{makeglossaries}. Los fragmentos de rúbrica y secciones compartidas no se compilan por separado.

También puedes trabajar en línea. En \texttt{Overleaf}, abre \texttt{Settings} > \texttt{Compiler} > \texttt{Main document}, selecciona el fichero \texttt{.tex} del formulario dentro de \fichero{PapeleoTFG/} y pulsa \texttt{Recompile}. Conserva el proyecto completo, incluido \fichero{Config/}, y revisa las condiciones de privacidad del servicio antes de subir datos personales o información confidencial. \fichero{00-README.pdf} explica cómo importar la distribución en esta plataforma.

Si utilizas \texttt{make}, ejecuta estas órdenes desde \fichero{PapeleoTFG/}. Necesitas \texttt{latexmk} y su intérprete \texttt{perl}, además de \texttt{make}, \texttt{bash} y la distribución de \LaTeX{}:

\medskip
\begin{minipage}{\linewidth}
\begin{verbatim}
make                         # Todos los formularios
make autorizaciones          # Publicación y confidencialidad
make convenios               # Convenio de cooperación
make evaluacion              # Informes, acta y rúbricas
make TFG-AnexoII-InformeTutorPlagio-UAH.pdf
make help                    # Opciones disponibles
\end{verbatim}
\end{minipage}
Cada fichero independiente genera un PDF con su mismo nombre. \texttt{make clean} elimina auxiliares y conserva los PDF; \texttt{make distclean} elimina también los PDF de formularios. La distribución ZIP/TGZ incluye esta guía solo en PDF y la conserva al limpiar. Puedes consultarla directamente, sin compilar nada.

\section{Antes de presentar los documentos}

Dedica unos minutos a revisar el resultado, aunque se haya generado sin errores:
\begin{itemize}
  \item Comprueba que el formulario y su versión corresponden a tu titulación y al procedimiento que te hayan indicado. No todos los documentos se utilizan en todos los trabajos.
  \item Revisa nombres, DNI, afiliaciones, títulos, fechas, géneros, cotutor y miembros del tribunal. No dejes los datos de ejemplo de la plantilla.
  \item Completa campos vacíos y \hl{N/D}, análisis de plagio, notas, observaciones y decisiones de publicación. Revisa el género también en los tratamientos escritos dentro de nombres o cargos, que no se reescriben automáticamente.
  \item Obtén las firmas que correspondan y sigue el procedimiento de entrega que te indique la EPS-UAH. La plantilla genera documentos, no los firma ni los presenta por ti.
\end{itemize}
\end{minipage}
\end{document}
